% chapters/05/section-01.tex
\section{Posisi dan Kecepatan sebagai Vektor}
\label{sec:posisi-vektor}

Bab~II telah memperkenalkan posisi, kecepatan, dan percepatan sebagai besaran
bertingkat yang saling berkaitan. Bagian ini mengubah gagasan tersebut menjadi
jenis data yang dapat dipakai di dalam program.

\subsection{Mengapa Pasangan Angka Perlu Disatukan}

Sebuah objek pada bidang datar memiliki posisi yang dinyatakan oleh dua angka:
jarak mendatar dan jarak tegak. Cara paling sederhana adalah menyimpannya
pada dua variabel terpisah.

\begin{kodekeluaran}
position_x = 200.0
position_y = 130.0
velocity_x = 200.0
velocity_y = 0.0
\end{kodekeluaran}

Cara itu dapat dipakai, tetapi cepat menjadi merepotkan. Setiap operasi harus
ditulis dua kali, dan setiap kali sebuah fungsi menerima posisi, ia harus
menerima dua argumen. Ketika jumlah objek bertambah, jumlah variabel bertambah
dua kali lipat, dan kekeliruan menukar pasangan angka---misalnya menambahkan
kecepatan mendatar ke posisi tegak---menjadi sulit terlihat karena keduanya
bertipe sama.

Pygame menyediakan jenis data \code{pygame.Vector2} yang menyatukan kedua
angka tersebut menjadi satu nilai. Posisi dan kecepatan pada
Gambar~\ref{fig:vektor-kecepatan} dinyatakan sebagai vektor, dan seluruh
perhitungan dilakukan atas vektor, bukan atas komponennya satu per satu.

\begin{figure}[htbp]
\centering
% chapters/05/figures/vektor-kecepatan.tikz
% Posisi sebagai vektor, kecepatan sebagai vektor, dan perpindahan yang
% dihasilkan selama satu selang waktu. Perpindahan digambar segaris dengan
% kecepatan karena nilainya memang kecepatan dikalikan selang waktu.
\begin{tikzpicture}[
  panahposisi/.style={-{Stealth[length=2.4mm]}, very thick, draw=ubPrimary},
  panahkecepatan/.style={-{Stealth[length=2.4mm]}, very thick, draw=ubAccent},
  panahpindah/.style={-{Stealth[length=2.2mm]}, thick, draw=ubGreen, dashed},
  kisi/.style={draw=ubLine!50, very thin},
  label/.style={font=\scriptsize, align=left}
]

% Kisi latar sebagai gambaran koordinat layar.
\foreach \x in {0,1,...,10} \draw[kisi] (\x,0) -- (\x,-4);
\foreach \y in {0,-1,...,-4} \draw[kisi] (0,\y) -- (10,\y);

% Sumbu koordinat layar: x ke kanan, y ke bawah.
\draw[-{Stealth[length=2mm]}, thick, draw=ubLine] (0,0) -- (10.4,0)
  node[anchor=north west, font=\scriptsize, ubLine] {$x$};
\draw[-{Stealth[length=2mm]}, thick, draw=ubLine] (0,0) -- (0,-4.3)
  node[anchor=north east, font=\scriptsize, ubLine] {$y$};

% Titik-titik: posisi awal p, posisi sesudah satu selang waktu p', dan ujung
% vektor kecepatan. Perpindahan p ke p' segaris dengan kecepatan p ke v.
\coordinate (p) at (2,-2.6);
\coordinate (pp) at (3.36,-2.04);
\coordinate (v) at (5.4,-1.2);

\draw[panahposisi] (0,0) -- (p);
\fill[ubPrimary] (p) circle (1.8pt);
\node[label, ubPrimary, anchor=south east] at (1.9,-2.75)
  {posisi $p = (200;\ 130)$};

\draw[panahpindah] (p) -- (pp);
\fill[ubGreen] (pp) circle (1.6pt);

\draw[panahkecepatan] (p) -- (v);
\node[label, ubAccent, anchor=south west] at (4.2,-1.45)
  {kecepatan $v$ dalam piksel per detik};

\node[label, ubGreen, anchor=north west] at (2.35,-2.95)
  {perpindahan $v \cdot \Delta t$};

\node[label, ubGreen, anchor=west] at (3.5,-2.0) {$p'$};
\draw[ubGreen, very thin] (pp) -- (3.5,-2.0);
\end{tikzpicture}

\caption{Posisi, kecepatan, dan perpindahan sebagai vektor}
\label{fig:vektor-kecepatan}
\end{figure}

\subsection{Operasi Vektor yang Sering Dipakai}

Tabel~\ref{tab:operasi-vector2} memuat operasi vektor yang akan Anda pakai
sepanjang Blok~I dan Blok~II. Seluruh operasi itu disediakan oleh Pygame,
sehingga tidak perlu ditulis sendiri \cite{pygame-docs}.

\begin{table}[htbp]
\centering
\caption{Operasi \code{pygame.Vector2} yang paling sering dipakai}
\label{tab:operasi-vector2}
{\setlength{\tabcolsep}{4pt}\footnotesize%
\begin{tabularx}{\textwidth}{@{}L{5.2cm} Y@{}}
\toprule
\headrow \textbf{Penulisan} & \textbf{Arti dan contoh} \\
\midrule
\code{pygame.Vector2(200, 130)} &
Membuat vektor baru. Contoh di atas berarti 200 satuan mendatar dan 130 satuan
tegak dari titik asal. \\
\addlinespace
\code{position.x}, \code{position.y} &
Membaca atau mengubah salah satu komponen. Berguna untuk membatasi posisi pada
tepi jendela. \\
\addlinespace
\code{position + velocity * dt} &
Menjumlahkan vektor dan mengalikannya dengan satu angka. Perhitungan dilakukan
per komponen secara otomatis. \\
\addlinespace
\code{velocity.length()} &
Panjang vektor, yaitu besarnya kecepatan dalam piksel per detik. \\
\addlinespace
\code{velocity.length_squared()} &
Kuadrat panjang vektor. Dipakai untuk memeriksa apakah vektor bernilai nol,
tanpa menghitung akar. \\
\addlinespace
\code{direction.normalize()} &
Vektor baru yang arahnya sama tetapi panjangnya satu. Vektor asalnya tidak
berubah. \\
\addlinespace
\code{a.distance_to(b)} &
Jarak antara dua titik. Berguna nanti untuk mengukur jarak pemain terhadap
objek lain. \\
\bottomrule
\end{tabularx}}
\end{table}

Operasi penjumlahan dan perkalian dengan satu angka bekerja per komponen. Jadi
\code{pygame.Vector2(200, 130) + pygame.Vector2(10, -5)} menghasilkan vektor
bernilai (210, 125), dan \code{pygame.Vector2(10, 5) * 3} menghasilkan vektor
bernilai (30, 15). Cara berpikir yang perlu dibangun adalah memperlakukan
vektor sebagai satu nilai, bukan sebagai dua angka yang kebetulan ditulis
berdampingan.

\subsection{Satu Kekhasan Koordinat Layar}

Satu hal yang sudah disebut pada Bagian~2.3 perlu dikembalikan ke hadapan Anda
sekarang, karena akibatnya akan terasa pada kode.

Pada koordinat layar, sumbu tegak menuju ke \textbf{bawah}. Karena itu, gerak
ke atas dinyatakan sebagai \textbf{pengurangan} nilai tegak, bukan
penambahan.

\begin{kodekeluaran}
if keys[pygame.K_UP] or keys[pygame.K_w]:
    direction.y -= 1        # ke atas, bukan ke bawah
\end{kodekeluaran}

Kekeliruan pada baris itu merupakan salah satu kesalahan pertama yang paling
sering terjadi: tombol atas ditekan, tetapi objek bergerak ke bawah. Apabila
Anda mengalaminya, periksa lebih dahulu tanda pada komponen tegak, sebelum
mencari penyebab lain yang lebih rumit.

\begin{catatabox}
Fungsi \code{normalize()} mengembalikan vektor \textbf{baru} dan
tidak mengubah vektor asalnya. Menulis \code{direction.normalize()} tanpa
menyimpan hasilnya tidak berpengaruh apa pun. Apabila Anda memang ingin
mengubah vektor asalnya, gunakan \code{normalize_ip()} yang bekerja di tempat.
Perbedaan keduanya mudah terlewat dan pernah menjadi sumber kebingungan.
\end{catatabox}
